\scanpage{25}
1) gr. alg $\boxed{G}$ sur $\mathbb Q_p$

2) hom continu \(\boxed{\rho:\pi\unsure{\mathbin{-}|\longrightarrow}G(\mathbb Q_p)}\)

d’où $G_\Omega^{!}\mathbin{-}\unsure{Q|\Phi|\mathfrak{P}}_\Omega^{!}\mathbin{-}G_\Omega$

% slanted
\reversemarginpar
\margin{unique\,$\left\{\vphantom{\begin{matrix}1\\1\\1\\1\end{matrix}}\right.$}
\normalmarginpar
3) \(\boxed{G_{m_\Omega}\xrightarrow{\unsure{j|J}^{!}}G_\Omega^{!}}\) \qquad d’où
\diagram{\begin{tikzcd}[ampersand replacement=\&,row sep=3pt,column sep=5pt]
\&Q_\Omega=\unsure{\mathcal T|\tau}\mathbin{\wedge}_{G_{m_\Omega}}Q_\Omega^{!}\arrow[dl,dash]\arrow[dr,dash]\&\\
G_\Omega^{\prime\prime}=\mathcal T\bigl(\mathbin{\wedge}^{\ast}_{G_{m_\Omega}}(G_\Omega^{!},\unsure{i^{-1}|\mathrm{int}})\bigr)\&\&G_\Omega\\
j_\Omega:G_m\longrightarrow G_\Omega^{\prime\prime}\&\&
\end{tikzcd}}
\deleted{4)}

4) \(\boxed{Q_L\ \text{sous }G_L}\) descendant $Q_\Omega$ sous $G_\Omega$
\[
U_L^{\prime\prime}\subset G_L^{\prime\prime}\mathbin{-}Q_L\mathbin{-}G_L,\quad
j_L:G_{m_L}\unsure{\longrightarrow|\mathbin{-}}G_L^{\prime\prime}
\]

5) \deleted{\unsure{Un ensemble} $\boxed{\Delta}$ d’opérations \inserted{$\rho$} de $\Gamma$ $(\supset\pi)$ sur $Q_\Omega$ \unsure{compatibles} \illegible{2} prolongeant celle de $\pi$, \inserted{d’où $\Gamma\longrightarrow G_L^{\prime\prime}(\unsure{\Omega|L})=G_\Omega^{\prime\prime}(\Omega)$} comp. à \unsure{une action|son action|celle} de $\Gamma$ sur $\Omega$.

tel que si $\rho$ est stable par \unsure{les|les autres} formant orbite sous $U_L^{\prime\prime}(L)$ opérant par (pour $u\in U_L^{\prime\prime}(L)$)
\[
(u.\rho)(\unsure{\gamma})=\tau(u)\bigl(\rho(\gamma)\bigr)\tau(u)^{-1}=(\unsure{\operatorname{int}\tau(u)}).\rho(\gamma)
\]
i.e. $u\unsure{\mapsto|\mathbin{-}|\mathbin{.}}\rho=\operatorname{int}\tau(u)\rho)$

[ NB si $u.\rho=\rho$ i.e. $\tau(u),\rho$ commutent : \unsure{$\tau(u)$} pour \unsure{tout} ]}

Un torseur (ensembliste) \inserted{à droite} $R_0\inserted{\Delta}$ sous $U_L^{\prime\prime}(L)$, et un $U_L^{\prime\prime}(L)$-\uline{hom de}
\[
\left[\begin{aligned}
\rho:\Delta\longrightarrow{}&\text{Op. de }\Gamma\text{ sur }(\Omega,Q_L,G_L)\\
&\text{qui sont compatibles avec}\\
&\text{\deleted{\unsure{l’}}l’op qu’on a de }\Gamma\text{ sur }\Omega,\\
&\text{et celle de }\pi\unsure{\subset|\in}\Gamma\text{ sur }Q_L\ldots
\end{aligned}\right.
\]
lorsque on fait opérer $U_L^{\prime\prime}(L)$ \inserted{à droite} sur le deuxième ens. par
\[
u.\alpha=\tau(\alpha)^{-1}u\tau(\alpha)^{\deleted{\unsure{-1}}}
\]
% archive sheet number: 24
